\documentclass[11pt,a4paper]{article}
\usepackage[utf8]{inputenc}
\usepackage[vietnamese]{babel}
\usepackage{amsmath,amssymb}
\usepackage[version=4]{mhchem}
\usepackage{geometry}
\geometry{top=1.5cm,bottom=1.2cm,left=1.4cm,right=1.4cm,headheight=28pt,headsep=12pt,footskip=18pt}
\usepackage{graphicx}
\usepackage{xcolor}
\usepackage{tabularx}
\usepackage{array}
\usepackage{tikz}
\usepackage{fancyhdr}
\usepackage{enumitem}
\usepackage{multido}
\usepackage{tcolorbox}
\usepackage{colortbl}

\definecolor{hfRed}{HTML}{D32F2F}
\definecolor{hfGreen}{HTML}{8BC390}
\definecolor{mainGreen}{HTML}{2E7D32}
\definecolor{solBlue}{HTML}{1565C0}
\definecolor{tableGreen}{HTML}{E8F5E9}

\usepackage{ex_test}

\pagestyle{fancy}
\fancyhf{}
\renewcommand{\headrulewidth}{0.75pt}
\renewcommand{\headrule}{\hbox to\headwidth{\color{hfGreen}\leaders\hrule height \headrulewidth\hfill}}
\renewcommand{\footrulewidth}{0.75pt}
\renewcommand{\footrule}{\hbox to\headwidth{\color{hfGreen}\leaders\hrule height \footrulewidth\hfill}}

\fancyhead[L]{\small\itshape\color{hfRed}Tài liệu lưu hành nội bộ}
\fancyhead[R]{\small\color{hfRed}Thầy \textbf{TRẦN VĂN THẠNH} - 0777.470.803}

\fancyfoot[L]{\footnotesize\itshape\color{hfRed}CS1: 6/15 Nguyễn Hoàng, P. Kim Long, TP Huế \quad|\quad CS2: 24 Đặng Thái Thân, TP Huế}
\fancyfoot[R]{\small\textbf{\color{hfRed}Trang \thepage}}

\setlength{\parindent}{0pt}
\setlength{\parskip}{2pt}

\begin{document}

\begin{minipage}[t]{0.48\textwidth}
\textbf{SỞ GIÁO DỤC VÀ ĐÀO TẠO TP. HUẾ}\\
\textbf{TRƯỜNG THPT NGUYỄN CHÍ THANH}\\
\textit{ĐỀ CHÍNH THỨC}\\
\textit{(Đề thi có 04 trang)}
\end{minipage}
\hfill
\begin{minipage}[t]{0.48\textwidth}
\raggedleft
\textbf{\color{hfRed}ĐỀ KIỂM TRA GIỮA HỌC KỲ I – NĂM HỌC 2026 – 2027}\\
\textbf{MÔN: HOÁ HỌC. LỚP 11}\\
\textit{Thời gian làm bài: 45 phút (không kể thời gian phát đề)}\\
\textbf{Mã đề: 1002}
\end{minipage}

\vspace{4pt}
\noindent
\textit{Họ, tên thí sinh:} \dotfill \textit{Lớp:} \dotfill \textit{Số báo danh:} \dotfill

\vspace{2pt}
\noindent
\textit{\footnotesize (Cho nguyên tử khối của: $\ce{H = 1; C = 12; N = 14; O = 16; Na = 23; Al = 27; S = 32; Cl = 35{,}5; Zn = 65; Ba = 137}$).}

\vspace{2pt}
\hrule color{hfGreen} height 0.5pt

\vspace{4pt}
\noindent
\textbf{\large\color{mainGreen}PHẦN I. Câu trắc nghiệm nhiều phương án lựa chọn: (16 câu - 4 điểm)}

\textit{Thí sinh trả lời từ câu 1 đến câu 16. Mỗi câu thí sinh chỉ chọn một phương án.}
\vspace{4pt}

\begin{ex}%[1002-C1]
Hiện tượng nào sau đây không phải là tác hại chính của mưa acid?
\choice
{\True Làm gia tăng hiện tượng nóng lên toàn cầu}
{Làm hủy hoại rừng và thảm thực vật cây cối}
{Gây ăn mòn và hủy hoại các công trình kiến trúc}
{Gây hại và làm suy giảm sinh vật dưới nước}

\end{ex}

\begin{ex}%[1002-C2]
pH của dung dịch nào sau đây có giá trị nhỏ nhất?
\choice
{\True Dung dịch $\ce{HCl}$ $0{,}1\text{M}$.}
{Dung dịch $\ce{CH3COOH}$ $0{,}1\text{M}$.}
{Dung dịch $\ce{NaCl}$ $0{,}1\text{M}$.}
{Dung dịch $\ce{NaOH}$ $0{,}01\text{M}$.}

\end{ex}

\begin{ex}%[1002-C3]
Chuẩn độ dung dịch $\ce{NaOH}$ chưa biết chính xác nồng độ (biết nồng độ trong khoảng gần với $0{,}1\,\text{M}$) bằng dung dịch chuẩn $\ce{HCl}$ $0{,}1\,\text{M}$ với chỉ thị phenolphthalein. Phương trình thể hiện bản chất của phản ứng của quá trình chuẩn độ là ?
\choice
{$\ce{NaCl -> Na+ + Cl-}$}
{$\ce{Na+ + Cl- -> NaCl}$}
{$\ce{H2O <=> H+ + OH-}$}
{\True $\ce{H+ + OH- -> H2O}$}

\end{ex}

\begin{ex}%[1002-C4]
Cho $m$ gam $\ce{Al}$ phản ứng hoàn toàn với dung dịch $\ce{HNO3}$ loãng (dư), thu được $4{,}958$ lít khí $\ce{NO}$ (đkc, sản phẩm khử duy nhất). Giá trị của $m$ là
\choice
{4,5.}
{8,1}
{\True 5,4.}
{2,7}

\end{ex}

\begin{ex}%[1002-C5]
Nitrogen ($\ce{N2}$) thể hiện tính khử trong phản ứng với chất nào trong các chất sau?
\choice
{$\ce{H2}$.}
{$\ce{Li}$.}
{\True $\ce{O2}$.}
{$\ce{Mg}$}

\end{ex}

\begin{ex}%[1002-C6]
Tính chất hóa học cơ bản của ammonia ($\ce{NH3}$) là
\choice
{tính base mạnh, tính khử.}
{tính base yếu, tính oxi hóa.}
{tính base mạnh, tính oxi hóa.}
{\True tính khử mạnh, tính base yếu.}

\end{ex}

\begin{ex}%[1002-C7]
Các yếu tố có thể ảnh hưởng đến sự chuyển dịch cân bằng hoá học là :
\choice
{nồng độ, áp suất và diện tích bề mặt.}
{\True nồng độ, nhiệt độ và áp suất.}
{áp suất, nhiệt độ và chất xúc tác.}
{nồng độ, nhiệt độ và chất xúc tác.}

\end{ex}

\begin{ex}%[1002-C8]
Số oxi hoá của nitrogen ($\ce{N}$) trong phân tử nitric acid ($\ce{HNO3}$) là
\choice
{$+4$.}
{\True $+5$.}
{$+2$.}
{$+3$.}

\end{ex}

\begin{ex}%[1002-C9]
Hành động nào sau đây của con người không phải là nguyên nhân chính gây ra hiện tượng phú dưỡng?
\choice
{Xả thải nước thải chăn nuôi, công nghiệp chứa hàm lượng cao các nguyên tố dinh dưỡng ($\ce{N, P}$).}
{\True Đánh bắt thuỷ hải sản bằng chích điện và chất nổ.}
{Sử dụng dư thừa lượng phân bón chứa ion $\ce{PO4^{3-}}$ ngấm vào nguồn nước mặt.}
{Sử dụng dư thừa lượng phân bón chứa các ion $\ce{NO2^-}, \ce{NO3^-}$ trong sản xuất nông nghiệp.}

\end{ex}

\begin{ex}%[1002-C10]
Phát biểu nào sau đây về acid – base theo thuyết Bronsted – Lowry là đúng?
\choice
{\True Acid là chất cho proton ($\ce{H+}$).}
{Base là chất nhận electron.}
{Base là chất cho proton ($\ce{H+}$).}
{Acid là chất nhận proton ($\ce{H+}$).}

\end{ex}

\begin{ex}%[1002-C11]
Cho dung dịch nitric acid ($\ce{HNO3}$) tác dụng với các chất sau: $\ce{Cu, NaOH, C, FeO, Na2CO3}$. Khi phản ứng xảy ra, có bao nhiêu phản ứng trong đó nitric acid ($\ce{HNO3}$) đóng vai trò oxi hóa ?
\choice
{4.}
{5.}
{2.}
{\True 3.}

\end{ex}

\begin{ex}%[1002-C12]
Các dung dịch của muối, acid, base dẫn điện là do
\choice
{phân tử của chúng dẫn được điện.}
{có sự di chuyển của electron tạo thành dòng electron.}
{các ion hợp phần có khả năng dẫn điện.}
{\True muối, acid, base có khả năng phân li ra ion khi tan trong nước.}

\end{ex}

\begin{ex}%[1002-C13]
Hệ phản ứng sau ở trạng thái cân bằng : $\ce{H2(g) + I2(g) <=> 2HI(g)}$\\
Biểu thức hằng số cân bằng của phản ứng trên là
\choice
{$K_\text{c} = \dfrac{[2\ce{HI}]}{[\ce{H2}]\cdot[\ce{I2}]}$}
{$K_\text{c} = \dfrac{[\ce{H2}]\cdot[\ce{I2}]}{2[\ce{HI}]}$}
{$K_\text{c} = \dfrac{[\ce{H2}]\cdot[\ce{I2}]}{[\ce{HI}]^2}$}
{\True $K_\text{c} = \dfrac{[\ce{HI}]^2}{[\ce{H2}]\cdot[\ce{I2}]}$}

\end{ex}

\begin{ex}%[1002-C14]
Ứng dụng phổ biến và quan trọng nhất của muối ammonium trong đời sống và sản xuất là
\choice
{sản xuất nhiên liệu xăng sinh học và dầu diesel.}
{sản xuất xi măng và vật liệu chịu lửa trong xây dựng.}
{\True sản xuất phân bón hoá học, cung cấp nguyên tố nitrogen dinh dưỡng cho cây trồng.}
{điều chế kim loại nguyên chất bằng phương pháp nhiệt luyện.}

\end{ex}

\begin{ex}%[1002-C15]
Bóng cười (hay "funky ball") là một quả bóng được bơm đầy khí X. Khí X có vị ngọt nhẹ, không màu, gây tác động hưng phấn, ảo giác cười cho người hít phải. Lạm dụng bóng cười gây tác hại nghiêm trọng đến hệ thần kinh, tim mạch, tụt huyết áp, co giật và nguy cơ tử vong. Công thức hoá học của khí X là
\choice
{$\ce{O2}$.}
{$\ce{N2}$.}
{$\ce{NO}$.}
{\True $\ce{N2O}$.}

\end{ex}

\begin{ex}%[1002-C16]
Phản ứng nào sau đây là phản ứng thuận nghịch ?
\choice
{$\ce{Fe + 2HCl -> FeCl2 + H2}$.}
{$\ce{H2 + Cl2 -> 2HCl}$.}
{$\ce{2H2 + O2 -> 2H2O}$.}
{\True $\ce{N2 + 3H2 <=> 2NH3}$.}

\end{ex}

\vspace{6pt}
\noindent
\textbf{\large\color{mainGreen}PHẦN II. Câu trắc nghiệm đúng sai: (2 câu - 2 điểm)}

\textit{Thí sinh trả lời từ câu 1 đến câu 2. Trong mỗi ý a), b), c), d) ở mỗi câu thí sinh chọn đúng hoặc sai.}
\vspace{4pt}

\begin{ex}%[1002-P2-C1]
Trong dung dịch muối $\ce{Fe^{3+}}$ tồn tại cân bằng hóa học sau:
$$\ce{Fe^{3+} + 3H2O <=> Fe(OH)3 v + 3H+}$$
\choiceTF
{\True Trong phòng thí nghiệm, để bảo quản dung dịch $\ce{Fe^{3+}}$, người ta thường thêm vào bình đựng vài giọt dung dịch acid $\ce{HCl}$ để cân bằng chuyển dịch theo chiều nghịch}
{\True Dựa vào thuyết acid – base của Brønsted – Lowry, $\ce{Fe^{3+}}$ đóng vai trò acid}
{Thêm vào dung dịch muối $\ce{Fe^{3+}}$ vài giọt dung dịch $\ce{NaOH}$ thì cân bằng chuyển dịch theo chiều nghịch}
{\True Các loại đất có chứa nhiều ion $\ce{Fe^{3+}}$ có giá trị pH thấp hay còn gọi là đất chua}

\end{ex}

\begin{ex}%[1002-P2-C2]
Trong những cơn mưa dông kèm sấm sét, sự chuyển hoá của nitrogen đơn chất xảy ra theo sơ đồ sau:
$$\ce{N2 ->[+O_2][(1)] NO ->[+O_2][(2)] NO2 ->[+O_2 + H_2O][(3)] HNO3 ->[(4)] H+ + NO3^-}$$
Ion nitrate ($\ce{NO3-}$) là một dạng đạm cây trồng hấp thụ được để sinh trưởng, phát triển.
\choiceTF
{\True Trong phản ứng (1) $\ce{N2}$ đóng vai trò là chất khử}
{Để tạo ra được $124\,\text{kg}$ ion nitrate cần dùng ở điều kiện chuẩn là $991{,}6\,\text{m}^3$ khí $\ce{N2}$. Biết hiệu suất của toàn bộ quá trình là $5\%$}
{\True Sơ đồ cũng áp dụng để giải thích sự tạo thành $\ce{HNO3}$ từ hoạt động đốt cháy nhiên liệu bằng không khí của phương tiện giao thông vận tải}
{Năng lượng trong sấm sét gây ra sự phân li của $\ce{HNO3}$ trong quá trình (4)}

\end{ex}

\vspace{6pt}
\noindent
\textbf{\large\color{mainGreen}PHẦN III. Câu trắc nghiệm trả lời ngắn: (4 câu - 1 điểm)}

\textit{Thí sinh trả lời từ câu 1 đến câu 4.}
\vspace{4pt}

\begin{ex}%[1002-P3-C1]
Trong công nghiệp, người ta sản xuất nitric acid ($\ce{HNO3}$) từ ammonia theo sơ đồ chuyển hoá sau:
$$\ce{NH3 ->[+O_2,\, t^\circ,\, \text{xt}] NO ->[+O_2] NO2 ->[+O_2 + H_2O] HNO3}$$
Để điều chế $8{,}4$ tấn dung dịch nitric acid có nồng độ $60\%$ cần dùng bao nhiêu tấn ammonia? Biết rằng hiệu suất của toàn bộ quá trình sản xuất nitric acid theo sơ đồ trên là $80\%$.

\textbf{Trả lời: } \dotfill

\end{ex}

\begin{ex}%[1002-P3-C2]
Hòa tan hết $0{,}8$ gam $\ce{NaOH}$ vào nước để được $2000\,\text{mL}$ dung dịch ở $25^\circ\text{C}$. pH của dung dịch thu được bằng mấy?

\textbf{Trả lời: } \dotfill

\end{ex}

\begin{ex}%[1002-P3-C3]
Cho các chất (phân tử và ion): $\ce{NaOH, H2SO4, NH3, CO3^{2-}, S^{2-}, NH4+}$. Theo quan điểm của Bronsted – Lowry có bao nhiêu chất (phân tử và ion) có thể là base?

\textbf{Trả lời: } \dotfill

\end{ex}

\begin{ex}%[1002-P3-C4]
Cho cân bằng hoá học sau: $\ce{2SO2(g) + O2(g) <=> 2SO3(g)};\ \Delta_\text{r}H^\circ_{298} < 0$\\
Người ta thực hiện các biện pháp sau:
\begin{enumerate}[label=(\arabic*)]
\item Giảm nhiệt độ,
\item Tăng áp suất chung của hệ phản ứng,
\item Dùng thêm chất xúc tác $\ce{V2O5}$,
\item Tăng nồng độ $\ce{SO3}$,
\item Tăng nồng độ $\ce{SO2}$.
\end{enumerate}
Có bao nhiêu biện pháp làm cân bằng trên chuyển dịch theo chiều thuận?

\textbf{Trả lời: } \dotfill

\end{ex}

\vspace{6pt}
\noindent
\textbf{\large\color{mainGreen}PHẦN IV. TỰ LUẬN: (3 câu - 3 điểm)}

\textit{Thí sinh trả lời từ câu 1 đến câu 3.}
\vspace{4pt}

\begin{ex}%[1002-TL1]
\textbf{(1,0 điểm)} Hoàn thành các phương trình phản ứng sau (ghi rõ điều kiện phản ứng)
\begin{enumerate}[label=\textbf{\alph*)}]
\item $\ce{N2 + O2 ->}$
\item $\ce{NH3 + HCl ->}$
\item $\ce{HNO3 \text{đặc} + FeO ->}$
\item Nhiệt phân $\ce{NH4HCO3 ->}$
\end{enumerate}

\noindent\textit{Bài làm:}
\multido{\i=1+1}{7}{\noindent\makebox[\linewidth]{\dotfill}\\[1.5pt]}

\end{ex}

\begin{ex}%[1002-TL2]
\textbf{(1,0 điểm)} Cho phản ứng tổng hợp ammonia ở trạng thái cân bằng tại $430^\circ\text{C}$.
$$\ce{N2(g) + 3H2(g) <=> 2NH3(g)} .$$
Ở trạng thái cân bằng, nồng độ của $\ce{NH3}, \ce{N2}$ và $\ce{H2}$ lần lượt là $0{,}30\,\text{M}$; $0{,}05\,\text{M}$ và $0{,}10\,\text{M}$.\\
Tính hằng số cân bằng của phản ứng ở nhiệt độ trên.

\noindent\textit{Bài làm:}
\multido{\i=1+1}{6}{\noindent\makebox[\linewidth]{\dotfill}\\[1.5pt]}

\end{ex}

\begin{ex}%[1002-TL3]
\textbf{(1,0 điểm)} Hydrochloric acid ($\ce{HCl}$) là một chất được phát hiện trong dịch vị dạ dày có nồng độ $10^{-4} - 10^{-3}\,\text{mol/L}$ và độ pH duy trì ở mức $3 - 4$ đối với người bình thường. Nếu thiếu acid trong dạ dày thì thức ăn không chuyển hóa được lâu dần gây suy nhược cơ thể, nếu dư acid lâu ngày thì nó sẽ phá hủy đường ruột gây viêm loét dạ dày. Khi cơ thể dư acid $\ce{HCl}$, người bệnh có thể sử dụng thuốc giảm đau dạ dày có tên gọi là Nabica (thành phần chính là $\ce{NaHCO3}$), được dùng để trung hòa bớt lượng acid dư trong dạ dày.\\
Giả sử dịch vị dạ dày người bệnh chứa $1{,}5$ lít dung dịch hỗn hợp thức ăn lỏng, trong đó chứa $0{,}09125$ gam $\ce{HCl}$. Nếu khả năng tiêu thụ $\ce{NaHCO3}$ của cơ thể người bệnh là $70\%$ thì khối lượng $\ce{NaHCO3}$ người đó cần đưa vào cơ thể để duy trì độ pH trong dạ dày ở mức 3 là bao nhiêu?

\noindent\textit{Bài làm:}
\multido{\i=1+1}{10}{\noindent\makebox[\linewidth]{\dotfill}\\[1.5pt]}

\end{ex}

\vspace{8pt}
\begin{center}
\textbf{---------------------- HẾT ----------------------}
\end{center}
\begin{center}
\textit{\small - Thí sinh không được sử dụng tài liệu.\\ - Cán bộ coi thi không giải thích gì thêm.}
\end{center}

\end{document}
